\section{Lie algebra for proofs}
\label{app:LieAlgebra}
In this section, we introduce further mathematical preliminaries
as necessary to provide self-contained proofs of our theorems.
We omit rigorous proofs for these observations,
and interested readers are directed to related textbooks and research papers \cite{iserles2000lie,lee2003smooth, petrogradsky2007wreath, tu2008manifold, rotman2009introduction, hall2013lie, robbin2022introduction,liu2022transformers}.
We open the section with the Lie group and Lie group action to motivate our Lie-algebraic approach.
\subsection{Lie groups and Lie algebras}
\label{app:subsubsec_Lieintro}
A set $G$ is a semigroup when $G$ is closed under an associative binary operation $\cdot: G \times G \to G$, usually called a "multiplication" operation.
If $\exists e \in G$ such that $e \cdot g = g \cdot e = g$ for $\forall g \in G$, $G$ is a monoid.
Moreover, $G$ is a group if for $\forall g \in G$ there exists an element $g^{-1} \in G$ such that $g^{-1} \cdot g = g \cdot g^{-1} = e$.
\paragraph{Lie group and Lie group action}
A Lie group $G$ is a smooth manifold equipped with a group structure under smooth multiplication.
For $x,y \in G$ let us denote the left translation $L_x(y) = xy$.
From the smoothness of the multiplication,
$L_x: G \to G$ is a diffeomorphism of the underlying manifold $\mcal{G}=G$.
In other words, translation can be thought of $G$ \emph{acting} on $\mcal{G}$.
The notion of Lie group \emph{action} generally follows:
\begin{definition}[Lie group action]
\label{app:def_LieGroupAction}
    Given a Lie group $G$ and a smooth manifold $\mcal{M}$,
    a (left) Lie group action is a smooth map
    $\Gamma: G \times \mcal{M} \to \mcal{M}$ that satisfies
    $\Gamma(e,h) = h$ and $\Gamma(xy,h) = \Gamma(x, \Gamma(y,h))$
    for an identity $e \in G$, any two elements $x, y \in G$,
    and a point on manifold $h \in \mcal{M}$.
\end{definition}
Following, we define the \emph{orbit} of $h$ denoted as $\text{Orb}_G(h) := \{\Gamma(g,h) \mid g \in G\}$.
\emph{Stabilizer}, or an isotropy subgroup, is a subgroup of $G$ where $\text{Stab}_G(h):= \{g\in G \mid \Gamma(g,h) = h \}$.

While \cref{app:def_LieGroupAction} straightforwardly describes how $\Gamma$ operates on $\mcal{M}$,
it is often advantageous to consider $\Gamma$ as inducing a diffeomorphism on $\mcal{M}$.
We can naturally introduce a Lie group homomorphism
$\Pi: G \to \text{Diff}(\mcal{M})$ such that $\Pi(g)=\Gamma(g,\cdot)$,
where $\text{Diff}(\mcal{M})$ denotes collection of diffeomorphisms of $\mcal{M}$.
\paragraph{Lie bracket and Lie algebra}
Now one may consider a smooth vector field $V: G \to TG$
and their collection $\fr{X}(G)$.
Given the differential of the left translation at point $y$,
\begin{equation*}
    (dL_x)_y: T_yG \to T_{L_x(y)}G = T_{xy}G,
\end{equation*}
the pushforward $(L_x)_*V \in \fr{X}(G)$ is well-defined \cite{tu2008manifold,lee2003smooth}.
Then, $V$ is said to be left-invariant when $(L_x)_*V = V$.
In other words, a left-invariant vector field $V$
satisfies $\bigl((L_x)_*V\bigr)_{xy} = (dL_x)_yV_y = V_{xy}$ for $\forall x, y \in G$
and completely determined by its evaluation at identity, as $V_x = V_{xe} = (dL_x)_e V_e$.
Thus, a collection of left-invariant vector fields, $\fr{X}_L(G)$, has a one-to-one correspondence with $T_e G$.

As a tangent vector is defined by how it differentiates smooth functions,
a smooth vector field on $\mcal{M}$
is commonly identified as the derivation $C^{\infty}(\mcal{M}) \to C^{\infty}(\mcal{M})$.
Following Leibniz rule,
the Lie bracket naturally arises to satisfy the derivation property \cite{tu2008manifold}.
For $f, g \in C^{\infty}(\mcal{M})$,
\begin{align*}
    XY(fg) &= X((Yf)g + fYg) \\
    &= (XYf)g + (Yf)(Xg) + (Xf)(Yg) + f(XYg) \\
    &\neq (XYf)g + f(XYg).
\end{align*}
Observe the extra term $(Yf)(Xg) + (Xf)(Yg)$ obstructs the derivation property.
One can subtract $YX(fg)$ to cancel out the extra term,
and naturally $(XY-YX)$ is again a derivation of $C^{\infty}(\mcal{M})$.
\begin{definition}[Lie bracket of vector fields]
    Given two smooth vector fields $X, Y: M \to TM$ on an open subset $M$ of a smooth manifold $\mcal{M}$,
    their Lie bracket is defined as following:
    \begin{equation}
    \label{app:eq_LieBracketDef}
        [X,Y] f = (XY - YX)f = X(Yf) - Y(Xf), \quad \forall f \in C^\infty(M).
    \end{equation}
\end{definition}
$\fr{X}_L(G)$ is closed under the Lie bracket.
Now we are ready to associate the Lie algebra $\fr{g}$ to the Lie group $G$:
\begin{definition}[Lie algebra associated to Lie group]
\label{app:LieAlgebra_Tangent}
    For a Lie group $G$, we identify the vector space of Lie algebra $\fr{g}:=T_eG$.
    For any $X \in \fr{g}$, assign a unique $X^L \in \fr{X}_L(G)$ that satisfies $X^L(y) = (dL_y)_eX$ for $y \in G$.
    We define the Lie bracket operation on $\fr{g}$ such that for $X,Y \in \fr{g}$,
    \begin{equation*}
        [X,Y] := [X^L,Y^L](e).
    \end{equation*}
    Then a map $\pi^L:\fr{g} \to \fr{X}_L(G)$, $\pi^L(X) = X^L$ is a Lie algebra isomorphism.
\end{definition}
Naturally, Lie algebra $\fr{g}$ can also be identified as a collection of the left-invariant, first-order differential operators on $G$.
Note the Lie bracket operation depends on the identification: with $\fr{g} \cong T_eG$, $[X,Y] \in T_eG$.
On the other hand, identifying $\fr{g} \cong \fr{X}_L(G)$ produces $[X,Y] \in \fr{X}_LG$.
The identification should be clear from the context.

We can analogously define the right translation on Lie group $R_x(y) = yx$.
The differential of the right translation at point y gives $(dR_x)_y: T_yG \to T_{R_x(y)}G = T_{yx}G$,
and we denote the collection of right-invariant vector fields as $\fr{X}_R(G)$.
Notationally, one can prove that $[X^L, Y^L](e) = -[X^R, Y^R](e)$ \cite{lee2003smooth}.
In that, $\pi^R: \fr{g} \to \fr{X}_RG$, $\pi^R(X) = X^R$ swaps the Lie bracket sign.

\paragraph{Bridge to geometry}
As a concrete example, consider the open set $M$ in \cref{app:eq_LieBracketDef} with local coordinate charts $x=(x^1, \cdots, x^n)$
around point $p \in M$.
Then we can rewrite $X \in \fr{X}(M)$ in a coordinate expression:
\begin{equation*}
    X = \sum_{a=1}^n X^a \partial_a,
    \quad
    X^a = X(x^a) \in C^\infty(M),
\end{equation*}
and similarly for $Y \in \fr{X}(M)$.
Given the local coordinate, Lie brackets get a tangible interpretation.
One can exhaust the vector field Lie brackets with coordinate,
\begin{align*}
    [X,Y] = \sum_{c=1}^n [X,Y]^c \partial_c,
    \quad
    [X,Y]^c = X(Y^c) - Y(X^c),
\end{align*}
where $[X,Y]^c$ is again a coordinate function of $[X,Y] \in \fr{X}(M)$.
Evaluated at $p$, $[X,Y]^c(p) = X_p(Y^c) - Y_p(X^c)$.
Then from the differential it is transparent:
\begin{equation}
\label{app:eq_LieBracket_Directional}
    [X,Y]^c(p) = X_p(Y^c) - Y_p(X^c) = (dY^c)_p(X_p) - (dX^c)_p(Y_p).
\end{equation}
One may interpret above as a directional derivative of a smooth coordinate function
along the tangent vector direction \cite{lee2003smooth}.

\subsection{Representation and infinitesimal action}
\label{app:subsubsec_repprelim}
\paragraph{Lie group representation}
The abstract notion provided at \cref{app:subsubsec_Lieintro} is often easier to work with a tangible object,
and matrix is one of such.
As $g, g^{-1} \in G$, let us consider the collection of any invertible matrices in $\bb{R}^{n \times n}$.
Over the real field, we name the collection as the general linear group and denotes as $\G{GL}{n;\bb{R}}$.
Matrix Lie group is defined as a closed subgroup of $\G{GL}{n;\bb{R}}$,
and many Lie groups in application arise as one.

These are the case of \emph{Lie group representations} \cite{lee2003smooth, hall2013lie}.
Given a real vector space $V$,
a representation of Lie group $G$ is a Lie group homomorphism $\Pi: G \to \G{GL}{V}$.
When $\dim V = n$, we can identify $\G{GL}{V}$ with $\G{GL}{n;\bb{R}}$ via basis choice.
Recall \cref{app:def_LieGroupAction} that representation is just a linear Lie group action of $G$ on $V$.
\paragraph{Lie algebra representation}
With \cref{app:LieAlgebra_Tangent} we can think of the corresponding \emph{Lie algebra representations}.
For a Lie group $G$ associated with $\fr{g}$,
consider $X \in \fr{g} \cong T_eG$ and corresponding $X^L \in \fr{X}_L(G)$.
Let $c_X: \bb{R} \to G$ be the integral curve of $X^L$ with $c_X(0) = e \in G$, starting at the identity.
Then we can construct the exponential map,
which is a canonical smooth map from the Lie algebra to its integrating group \cite{lee2003smooth}:
\begin{equation*}
    \exp: \fr{g} \to G, \quad \exp(X):= c_X(1),
\end{equation*}
and we write $\exp(tX) = c_X(t)$ so that $\frac{d}{dt}\vert_{t=0} \exp(tX) = X$.
Notice that $X$ corresponds to the initial velocity of $c_X(t)$.

One concrete example is the matrix Lie algebra $\fr{gl}(n,\bb{R})$, which is a collection of any matrices in $\bb{R}^{n \times n}$.
With the Lie group representation:
\begin{equation*}
    \Pi: G \to \G{GL}{n;\bb{R}}, \quad \Pi_*: T_eG \to T_{\Pi(e)}\G{GL}{n;\bb{R}} = T_{I_n}\G{GL}{n;\bb{R}},
\end{equation*}
where we identified $\G{GL}{n;\bb{R}}$-associated Lie algebra
$\fr{gl}(n,\bb{R}) \cong T_{I_n}\G{GL}{n;\bb{R}}$.
Then the Lie algebra representation is naturally induced:
\begin{equation*}
    \pi:= (d \Pi)_e: \fr{g} \to \fr{gl}(n,\bb{R}),
\end{equation*}
in that, $\pi$ is the differential of $\Pi$ at the identity.
For the matrix Lie algebra, exponential map is the matrix exponential $\text{expm}$:
\begin{equation}
\label{app:eq_matrixLie_exponential_action}
    \Pi(\exp(X)) = \text{expm} (\pi(X)), \: \pi(X) = \frac{d}{dt}\bigg\vert_{t=0}\Pi(\exp(tX)).
\end{equation}
We extend the differential of representation to the differential of action,
or \emph{infinitesimal} action, in the following paragraph.
\paragraph{Fundamental vector field}
Consider the Lie group $G$ acts on itself by left multiplication.
For $h \in G$ and $X \in \fr{g}$, let us define a vector field $X^\sharp$:
\begin{equation*}
    X^\sharp(h) = \frac{d}{dt} \bigg \vert_{t=0} \exp(-tX)h = -(dR_h)_eX \in T_hG.
\end{equation*}
Observe that $X^\sharp(h)$ recovers familiar ODE $\dot{h} = Xh$
up to sign change, when $G \subseteq \G{GL}{n,\bb{R}}$.

We can extend above to an arbitrary smooth finite manifold $\mcal{M}$,
replacing canonical left action on $G$ to a chosen left Lie group action $\Gamma: G \times \mcal{M} \to \mcal{M}$.
For each $X \in \fr{g}$ and $h \in \mcal{M}$, we define the fundamental vector field $X^\sharp \in \fr{X}(\mcal{M})$:
\begin{equation}
\label{app:eq_fundamental_vector_field}
    X^\sharp (h) := \frac{d}{dt}\bigg\vert_{t=0}\Gamma(\exp(-tX), h) \in T_h\mcal{M},
\end{equation}
then the map $\gamma: \fr{g} \to \fr{X}(\mcal{M})$,
$\gamma(X)= X^\sharp$ is a Lie algebra homomorphism \cite{iserles2000lie, robbin2022introduction}.
Therefore $\gamma$ is an \emph{infinitesimal} Lie group action, and considered as Lie algebra action on $\mcal{M}$.
\begin{remark}
\label{app:rem_Bracket_Sign}
Many text writes $X^\sharp$ in a \emph{right-invariant} vector field form \cite{robbin2022introduction}:
\begin{equation}
\label{app:eq_fundamental_vector_field2}
    X^\sharp (h) := \frac{d}{dt}\bigg\vert_{t=0}\Gamma(\exp(tX), h) \in T_h\mcal{M},
\end{equation}
and the map $\gamma(X)= X^\sharp$ is a Lie algebra \emph{anti}homomorphism \cite{iserles2000lie, robbin2022introduction}.
Under this definition, one would observe that $[\pi(X), \pi(Y)] = \pi([X,Y])$, a matrix commutator,
aligns with $[X^\sharp, Y^\sharp] = -[X,Y]^\sharp$ up to sign change.
In practice, the sign can be absorbed and need not be worried.
Some literature swap the Lie bracket direction to handle the sign change,
but this is just a choice of convention and does not obstruct any fundamental logic.
\end{remark}
\paragraph{Bridge to homogeneous ODE}
\label{app:bridge_to_hom_LPV}
As observed above, when a Lie group $G \subseteq \G{GL}{n,\bb{R}}$ left-acts on a smooth manifold $\mcal{M} \cong \bb{R}^n$,
the fundamental vector field recovers homogeneous ODE.
For a point $h \in \mcal{M}$ and $\mbf{A} \in \fr{g} \subseteq \fr{gl}(n,\bb{R})$, we select $\Gamma$ to be a left multiplication.
Then obviously:
\begin{equation*}
    \dot{h} = \frac{d}{ds}\bigg\vert_{s=0}\Gamma(\exp(s\mbf{A}), h) = \mbf{A}h,
\end{equation*}
which corresponds to the linear controlled homogeneous ODE, $\dot{h}(t) = \mbf{A}(x(t))h(t)$, given a fixed input $x(t)$.

Critically, when $G$ acts on itself, we can recover the state-transition matrix $\Phi$.
Recall \cref{def:Phi}, $\Phi$ satisfies:
\begin{equation}
\label{app:eq_homogeneous_matrix_state}
    \dot{\Phi}(t,0) = \mbf{A}(x(t))\Phi(t,0), \quad \Phi(0,0) = I_n,
\end{equation}
then \cref{app:eq_homogeneous_matrix_state} forms a homogeneous ODE system with $\Phi$ being its latent state.
This observation is very useful:
it is immediate that for any admissible input, $\Phi(t,0) \in G$.
In a sense, one may understand \cref{app:eq_homogeneous_matrix_state}
as tracking \emph{every} possible state evolution, rather than a single state.
This is a conventional perspective in geometric control theory,
commonly referred as control theory on Lie groups \cite{jurdjevic1972control, iserles2008magnus}.

\subsection{Lie algebra extension}
\label{app:subsubsec_extensionprelim}
Analogous to the notion of the group extension, Lie algebra extension formalizes \emph{gluing} two Lie algebras together.
Here we introduce the notion of Lie algebra extension \cite{weibel1994introduction,rotman2009introduction, PhysicsExtension, beckett2022symplectic}.
\begin{definition}[Lie algebra extension]
    \label{app:def_LieAlgebraExtension}
    The Lie algebra $\tilde{\fr{h}}$ is an extension of base Lie algebra $\fr{h}$ by (finite-dimensional) kernel Lie algebra $\fr{a}$ if there exists a short exact sequence of Lie algebras:
    \begin{equation}
    \label{app:eq_LieAlgebraExtension_shortsequence}
        0 \to \mathfrak{a} \xrightarrow{i} \tilde{\mathfrak{h}} \xrightarrow{\theta} \mathfrak{h} \to 0.
    \end{equation}
\end{definition}
Redundantly, $i: \fr{a} \to \tilde{\fr{h}}$ is a Lie algebra monomorphism,
while $\theta: \tilde{\fr{h}} \to \fr{h}$ is a Lie algebra epimorphism.
The short exact sequence satisfies $\text{im}\; i = \ker \theta$,
thus $i(\fr{a})$ is an ideal in $\tilde{\fr{h}}$ and $\tilde{\fr{h}}/i(\fr{a}) \cong \fr{h}$.
Following common convention, we identify $\fr{a}$ with $i(\fr{a})$ without confusion.
\begin{remark}[Tower of abelian extension]
\label{app:rem_solvable_k_tower}
A solvable Lie algebra with derived series $\fr{g} = \fr{g}^{(0)} \supset \fr{g}^{(1)} \supset \cdots \supset \fr{g}^{(k)}=\{0\}$
is an extension:
\begin{equation*}
    0 \to \fr{g}^{(n)} \to \fr{g} \to \fr{g} / \fr{g}^{(n)} \to 0.
\end{equation*}
When $n = k-1$, $\fr{g}^{(k-1)}$ is an abelian Lie (sub-)algebra and
$\fr{g} / \fr{g}^{(k-1)}$ is a solvable Lie algebra of derived length $k-1$.
In iteration, we recover a $k$-length \emph{tower} of abelian Lie algebra extension.
\paragraph{Split extension}
Understanding how $\fr{h}$ acts on $\fr{a}$ is critical in characterizing $\tilde{\fr{h}}$.
We begin with a simple case.
\end{remark}
\begin{definition}[Split extension]
    In \cref{app:eq_LieAlgebraExtension_shortsequence}
    consider a linear section $s: \fr{h} \to \tilde{\fr{h}}$ that satisfies $\theta \circ s = 1_\fr{h}$, an identity map.
    If $s$ is a Lie algebra homomorphism, \cref{app:eq_LieAlgebraExtension_shortsequence} is \emph{split}.
\end{definition}
If Lie algebra extension splits, then
we say $\tilde{\fr{h}}$ is a \emph{semidirect product} (or sum) of $\fr{h}$ by $\fr{a}$
and we denote as $\tilde{\fr{h}} = \fr{h} \ltimes \fr{a}$
(or $\fr{a} \rtimes \fr{h}$) \cite{rotman2009introduction}.\footnote{Note
that the convention of sign direction is not accordant across references \cite{petrogradsky2007wreath, rotman2009introduction}.
here we emphasize that $\fr{h}$ acts on $\fr{a}$.}

Consider $X,Y \in \fr{h}$ and $u,v \in \fr{a}$.
$\tilde{\fr{h}}$ acts on $\fr{a}$ by derivation, as $\fr{a}$ is an ideal in $\tilde{\fr{h}}$.
As the linear section $s$ is a Lie algebra homomorphism,
$s(X) \in \tilde{\fr{h}}$ naturally acts on $\fr{a}$ by derivation too.
Thus, the induced action of $\fr{h}$ on $\fr{a}$ is straightforward here:
\begin{equation*}
    \rho: \fr{h} \to \text{Der}(\fr{a}),
    \quad
    \rho(X)v = [s(X),v],
\end{equation*}
and Lie bracket on split extension follows:
\begin{align}
\label{app:eq_split_bracket}
    [(X, u), (Y, v)]_{\tilde{\fr{h}}} &= ([X,Y]_\fr{h}, [u,v]_\fr{a} + \rho(X)v - \rho(Y)u).
\end{align}
One special case of split extension is the trivial extension where $[(X,u),(Y,v)] = ([X,Y],[u,v])$,
and we name it as a \emph{direct product} (or sum) of $\fr{h}$ and $\fr{a}$.

\paragraph{Bridge to inhomogeneous ODE}
\label{app:bridge_to_general_LPV}
Similar to \cref{app:bridge_to_hom_LPV}, affine vector field naturally arises from the semidirect product.
Consider an inhomogeneous ODE $\dot{h}(t) = \mbf{A}(x(t))h(t) + \mbf{b}(x(t))$.
As noted often \cite{krener1975bilinear, krener1977decomposition, iserles2000lie},
one can augment the state to \emph{homogenize} the ODE:
\begin{equation*}
    \bar{h}(t) = \begin{pmatrix}
    h(t) \\
    1
\end{pmatrix},
\quad
\dot{\bar{h}}(t)
= \begin{pmatrix}
        \mbf{A}(x(t)) & \mbf{b}(x(t)) \\
        0 & 0
    \end{pmatrix}
    \:
    \bar{h}(t).
\end{equation*}
Then it is immediate that the generator class is affine.
Consider a finite-dimensional Lie group of affine transformation, $\G{Aff}{n, \bb{R}}$.
Its associated Lie algebra $\fr{aff}(n, \bb{R})$ is the canonical example of the split extension:
\begin{equation*}
    0 \to \bb{R}^n \to \fr{aff}(n, \bb{R}) \to \fr{gl}(n,\bb{R}) \to 0,
\end{equation*}
such that $(A,b) \in \fr{gl}(n,\bb{R}) \ltimes \bb{R}^n$.
With a standard matrix embedding,
\begin{equation*}
    (A,b) \longmapsto \begin{pmatrix}
        A & b \\
        0 & 0
    \end{pmatrix},
\end{equation*}
which corresponds to the generator class of the homogenized ODE above.

Let us confirm that realized SSM respects the split extension bracket rule in \cref{app:eq_split_bracket}.
From the matrix representation above, one can simply check their matrix commutator.
With two fixed inputs $x$, $y$:
\begin{equation*}
    \bar{\mbf{A}}(x) = \begin{pmatrix}
        \mbf{A}(x) & \mbf{b}(x) \\
        0 & 0
    \end{pmatrix},
    \quad
    [\bar{\mbf{A}}(x), \bar{\mbf{A}}(y)]
    =
    \begin{pmatrix}
        [\mbf{A}(x), \mbf{A}(y)]
        &
        \mbf{A}(x)\mbf{b}(y) - \mbf{A}(y)\mbf{b}(x) \\
        0 & 0
    \end{pmatrix},
\end{equation*}
especially when $\text{Lie}{\{\mbf{A}(x)\}}$ is abelian, $[\mbf{A}(x), \mbf{A}(y)] = 0$ and only the top right entry is nonzero.

Now consider vector field representation. For the brevity, let us keep $\text{Lie}{\{\mbf{A}(x)\}}$ be abelian.
Recall \cref{app:eq_LieBracket_Directional}:
\begin{align}
X(h) &= \mbf{A}(x)h + \mbf{b}(x), \quad Y(h) = \mbf{A}(y)h + \mbf{b}(y) \nonumber \\
    [X,Y]^c &= \sum_{a=1}^{n}
    \left(
    \mbf{b}^a(x)\mbf{A}^{ca}(y)
    - \mbf{b}^a(y)\mbf{A}^{ca}(x)
    \right) \nonumber \\ 
    [X,Y] &= \mbf{A}(y)\mbf{b}(x) - \mbf{A}(x)\mbf{b}(y),\label{app:eq_LPV_affinebracket}
\end{align}
with sign change noted in \cref{app:rem_Bracket_Sign}.

Taken together with \cref{app:bridge_to_hom_LPV},
switching the vector state space
with the corresponding Lie group, or flow space,
is often referred as \emph{lifting} \cite{krener1977decomposition}.
Except for minor details in implementation, lifting does not lose any generality
and we can \emph{project} back to the state space \cite{iserles2000lie, iserles2008magnus}.
\paragraph{Central extension and beyond}
If all extension belongs to the split extension class, our discussion would be much easier.
However note that the linear section $s: \fr{h} \to \tilde{\fr{h}}$ needs not be a Lie algebra homomorphism.
Among non-split extension, central extension is another key class of the Lie algebra extension.
\begin{definition}[Central extension]
    In \cref{app:eq_LieAlgebraExtension_shortsequence} let $\fr{a}$ be abelian.
    When $i(\fr{a})$ is in the center of $\tilde{\fr{h}}$ such that $[\tilde{X}, u] = 0$ for any $\tilde{X} \in \tilde{\fr{h}}$
    and $u \in i(\fr{a})$, we call it as central extension.
\end{definition}
As $s$ is not necessarily Lie algebra homomorphism, we introduce a correction term, called \emph{2-cocycle} \cite{weibel1994introduction}:
\begin{equation}
\label{app:eq_2cocycle}
    \omega: \fr{h} \times \fr{h} \to \fr{a}, \quad
    \omega(X,Y) := [s(X),s(Y)] - s([X,Y]_\fr{h}).
\end{equation}
Recall \cref{app:eq_split_bracket}.
As $\fr{a}$ belongs to the center, the induced action $\rho: \fr{h} \to \text{Der}(\fr{a})$ is trivial.
Thus, in central extension, the bracket rule simplifies to:
\begin{equation*}
\label{app:eq_central_bracket}
[(X,u),(Y,v)] = ([X,Y], \omega(X,Y)).
\end{equation*}
When $s$ is Lie algebra homomorphism, central extension breaks down to trivial extension.

Taken together, the Lie bracket on general Lie algebra extension follows:
\begin{equation*}
    [(X, u), (Y, v)]_{\tilde{\fr{h}}} = [[X,Y]_\fr{h}, [u,v]_\fr{a} + \rho(X)v - \rho(Y)u + \omega(X,Y)].
\end{equation*}
Characterization of general Lie algebra extension is nontrivial.
At least under Lie group setting, the Kaloujnine-Krasner universal embedding theorem states that
any Lie group extensions embeds into the \emph{wreath} product of Lie group \cite{wells1976some,deval2024universal}.
Fortunately, following studies have extended the universal embedding theorem
to Lie algebra \cite{petrogradsky2007wreath, deval2024universal}.
\section{Control theory basics}
\label{app:LPV}
\subsection{Abstract control systems}
\label{app:control_detail}
Here we discuss abstract notions of controlled dynamical systems \cite{iserles2000lie}.
On a real smooth finite-dimensional manifold $\mcal{M}$, we consider a time-dependent tangent vector field:
\begin{equation*}
    F: [0,\tau] \times \mcal{M} \times \mcal{X} \to T\mcal{M},
\end{equation*}
which is locally Lipschitz in $h \in \mcal{M}$ and measurable in time $t \in [0,\tau]$.  
In this context, we call $\mcal{M}$ the \emph{state space} and $\mcal{X}$ the \emph{input space}.
Then, with a fixed input path $x: [0,T] \to \mcal{X}$ we obtain the controlled ODE:
\begin{equation}
    \label{eq:abstractODE}
    \dot{h}(t) = F_x(t,h(t)), \; F_x(t,h(t)) := F(t,h(t),x(t)),
\end{equation}
assuming the path is locally bounded and integrable \cite{iserles2000lie}.
We call a tuple $S = (\mcal{M}, F,\mcal{X})$ a \emph{controlled dynamical system}.
Lastly, we say the system is initialized when $h(0)$ is fixed.
Restricting $\mcal{M}$ to be on a Euclidean state space realizes SSM.
\subsection{Simulation}
\label{app:sim_detail}
Here we formally introduce the notion of simulation \cite{sussmann1976existence,hermann1977nonlinear, krener1977decomposition}.
\begin{definition}[Simulation]
\label{def:SSM_sim_equivalence}
    Let two initialized $S_0$, $S_1$ share an input space.
    Consider a smooth map $\varphi: \mcal{M}_0 \to \mcal{M}_1$
    such that $\varphi(h_0(t)) = h_1(t)$ for any admissible inputs.
    If $\varphi$ is a surjective submersion, we say $S_0$ \emph{simulates} $S_1$.
    Moreover, if $\varphi$ is (locally) diffeomorphic, we say $S_0$ is equivalent to $S_1$.
\end{definition}
Note that if $S_0$ simulates $S_1$ and $\dim S_0 = \dim S_0$,
then $\varphi$ is diffeomorphic.
Therefore, equivalence implies simulation.

For SSMs, $\varphi$ is restricted to be full-rank linear map.
Following, diffeomorphic $\varphi$ is invertible linear map.
Arbitrary smooth $\varphi$ would typically leave the SSM model class \cite{krener1977decomposition,toth2010modeling}.
When restricted to SSMs, we define simulation error correspondingly.
\begin{definition}[SSM Simulation error on horizon]
\label{def:SSM_sim_error}
    Consider initialized $\mbf{S}_0$, $\mbf{S}_1$ with $n_0 = \dim \mbf{S}_0 \ge \dim \mbf{S}_1 = n_1$.
    Let $\mcal{X}_{ad}$ be shared admissible input path space, with
    \begin{equation*}
        U := \{x \in \mcal{X}_{ad} \mid x: [0,T] \to \mcal{X}, \|x\| \le M\}
    \end{equation*}
    be measure bounded paths over chosen interval $[0,T]$.
    Then, we define simulation error:
    \begin{equation}
    \label{eq:SSM_sim_error}
    \Delta_h := \inf_{\mbf{P} \in \bb{R}^{n_1 \times n_0}} \sup_{x \in U} \bigl\| \mbf{P}h_0(T) - h_1(T) \bigr\|,
    \end{equation}
    where $\mbf{P}$ is a full-rank linear map.
\end{definition}
Note that $\mbf{P}$ can be replaced with invertible linear map $\mbf{T}$
by state augmentation on $\mbf{S}_1$, which leads to the standard state-space equivalence relation
in linear parameter-varying (LPV) system theory \cite{toth2010modeling}.
\subsection{Minimality}
\label{app:min_detail}
To understand whether one system simulates the other, we need a notion of minimality.
With a concept of Lie group action, minimality comes in a straightforward manner \cite{jurdjevic1972control, krener1977decomposition}.

Consider $5$-tuple $(\mcal{M},F,\mcal{X}_{ad},\mcal{Y},g)$
and its evolution operator $\Psi$.
Consider a collection $S = \{\Psi_x\}$ for $x \in \mcal{X}_{ad}$,
then one can notice $S \subset \text{Diff}(\mcal{M})$ is a semigroup.
Although any $\Psi_x(t,0)$ represent as an invertible matrix,
the existence of admissible input for $\Psi_{x'} = \Psi_x^{-1}$ is not guaranteed.
With completion, we denote a group $G$ be the smallest subgroup of $\text{Diff}(\mcal{M})$ containing $S$.

For a chosen point $h \in \mcal{M}$, consider a collection $S(h) = \{\Psi(h): \Psi \in S\}$ and $G(h) = \{\Psi(h): \Psi \in G\}$.
It is akin to the orbit of Lie group action:
intuitively, $S(h)$ is the set of accessible points on $\mcal{M}$ under admissible inputs.
We say the system is (weakly) \emph{controllable} at point $h$ when $S(h) = \mcal{M}$.
When satisfied, the condition implies $G(h) = \mcal{M}$.

Now consider the output map.
Given any two points $h_1, h_2 \in \mcal{M}$, denote $y_1(t) = g(\Psi(t,0)(h_1))$ and $y_2(t)$ correspondingly.
Then $h_1$ and $h_2$ are said to be indistinguishable when $y_1(t) = y_2(t)$ for any admissible inputs.
The system is \emph{observable} when $h_1$ and $h_2$ being indistinguishable implies $h_1 = h_2$.
Collectively, we say the system $(\mcal{M},F,\mcal{X}_{ad},\mcal{Y},g)$ is \emph{minimal}
when both (weakly) controllable and observable \cite{krener1977decomposition}.
\begin{remark}
    Often, the attention is restricted to open subset $U \subseteq \mcal{M}$ however small.
    Without global Lie group action and orbits, one can characterize the system with Lie algebras.
    For example, weak controllability is replaced by \emph{Lie algebra accessibility rank condition} \cite{sontag2013mathematical}.
\end{remark}
One may easily notice the notion of minimality in LTI system is a special case of the above.
We can similarly induce notion of minimality for SSMs,
commonly studied as a (quasi-)LPV system \cite{toth2010modeling}.
\subsection{Magnus series and path signature}
\label{app:CDE}
The Magnus series provides Lie-algebraic expansion of the state-transition matrix.
As various evolution operators in physical applications are indeed a integral curve on a Lie group,
the Magnus expansion (or its truncation) serves as a useful tool in numerical approximation problems \cite{iserles2000lie}.
Given a system $\mbf{S}$, we explicitly write down a few first terms:
\begin{align*}
    \Phi(t,0) &= \exp \Omega(t),
    \quad
    \Omega(t) = \sum \Omega_i(t)\\
    \Omega_1(t) &= \int_0^t \mbf{A}(t_1) \; dt_1 \\
    \Omega_2(t) &= \frac{1}{2}\int_0^t \; dt_1 \int_0^{t_1} \; dt_2
    [\mbf{A}(t_1), \mbf{A}(t_2)] \\
    \Omega_3(t) &= \frac{1}{6}\int_0^t \; dt_1 \int_0^{t_1} \; dt_2 \int_0^{t_2} \; dt_3
    \Bigl(
    [\mbf{A}(t_1), [\mbf{A}(t_2), \mbf{A}(t_3)]] + [\mbf{A}(t_3), [\mbf{A}(t_2), \mbf{A}(t_1)]]
    \Bigr)
    \\
    \cdots
\end{align*}
Observe how the lower central series contributes to each order of the Magnus series.
Specifically, the $i$-th order Magnus term involves a linear combination of elements in the $(i-1)$-th member of the lower central series, $\fr{g}^{i-1}$.
Consequently, if $\fr{g}$ is a class $c$ nilpotent Lie algebra, then the Magnus expansion of $\mbf{S}_\fr{g}$ truncates at order $c$. As a special case, abelian $\fr{g}$ is the class $1$ nilpotent Lie algebra, all commutators of $\mbf{S}_\fr{g}$ vanish and we have truncation at order $1$.

Recent progress in continuous-time neural network modeling leverages
rough path theory and path signature \cite{muca2024theoretical, walker2024log, walker2025structured}.
Path signature is deeply related to the Magnus series \cite{chevyrev2024multiplicative,chevyrev2025primer}.
For linear CDE, the path signature can be expressed as
a tensor exponential of the Magnus expansion \cite{friz2022unified}.
In other words, the Magnus expansion is a representation of \emph{log-signature} \cite{morrill2021neural, walker2024log}
when input variation is bounded.
Log-signature underscores the algebraic property of input path,
while the Magnus expansion prioritizes state transition.
The study of the Magnus series is relatively mature and often provides both theoretical and computational benefits
when the path signature admits the Magnus expansion \cite{friz2022unified}.